% ----------------------------------------------------------------------
%  3.1.3. Силы межатомного взаимодействия; 3.1.4. Энергия связи
% ----------------------------------------------------------------------
\subsection{Общая характеристика сил межатомного взаимодействия}\label{ssec:forces}

Строго задачу о взаимодействии атомов нужно решать квантовомеханически,
приближённо решая уравнение Шрёдингера для всей системы ядер и
электронов. Но общий критерий прост:
\begin{glavnoe}
Если энергия сблизившихся атомов меньше суммы энергий изолированных
атомов, между ними преобладает притяжение, и образуется устойчивая
молекула (или кристалл). Если при сближении энергия растёт~—
преобладает отталкивание, и связи нет.
\end{glavnoe}

Какие силы могут связывать атомы? Сравним порядки величин энергий
взаимодействия двух атомов на расстоянии $r\approx0{,}25$~нм
(табл.~\ref{tab:forces}).

\begin{table}[ht]
\centering
\caption{Порядки энергий взаимодействия двух атомов на расстоянии $\sim0{,}25$~нм}
\label{tab:forces}
\small
\begin{tabular}{@{}llc@{}}
\toprule
Взаимодействие & Оценка & Энергия, эВ \\
\midrule
гравитационное & $Gm^2/r$ (тяжёлые атомы) & $\sim10^{-31}$ \\
магнитное & $\mu_0\mu_\text{Б}^2/(4\pi r^3)$ & $\sim10^{-5}$--$10^{-6}$ \\
электростатическое & $e^2/(4\pi\varepsilon_0r)$ & $\sim5$--$7$ \\
\bottomrule
\end{tabular}
\end{table}

Энергия связи атомов в кристаллах~— порядка электронвольт. Этому
соответствует только электростатическое взаимодействие. Значит,
\vazhno{все типы химической связи имеют электрическую природу}; связь
осуществляется внешними (\term{валентными}) электронами атомов. Типы
связи (п.~\ref{ssec:bond-types}) различаются тем, \emph{как} именно
распределены валентные электроны.

\subsection{Энергия связи кристаллов}\label{ssec:binding}

\paragraph{Парный потенциал}
Пусть частицы взаимодействуют центральными силами~— энергия пары зависит
только от расстояния $r$ между ними. Её удобно представить суммой
притяжения и отталкивания:
\begin{keyeq}\label{eq:pair-pot}
U(r)=U_{\text{пр}}+U_{\text{от}}=-\frac A{r^m}+\frac B{r^n},\qquad A,B>0 .
\end{keyeq}
Притяжение действует на больших расстояниях, отталкивание (перекрытие
электронных оболочек) резко нарастает на малых, поэтому $n>m$. Для
ионного взаимодействия $m=1$ (кулоновское притяжение), для
ван-дер-ваальсова $m=6$; отталкивание часто описывают показателем
$n\approx9$--$12$ (потенциал Леннард-Джонса: $m=6$, $n=12$).

\begin{figure}[ht]
\centering
\begin{tikzpicture}
\begin{axis}[width=0.66\textwidth, height=6cm, xmin=0.8, xmax=2.6, ymin=-1.3, ymax=1.6,
  axis lines=middle, xlabel={$r$}, ylabel={$U$}, xtick={1.1225}, xticklabels={$r_0$}, ytick={-1}, yticklabels={$U_0$},
  every axis x label/.style={at={(ticklabel* cs:1)}, anchor=west},
  every axis y label/.style={at={(ticklabel* cs:1)}, anchor=south}, clip=true]
  \addplot[fmRed, dashed, thick, domain=0.85:2.6, samples=100] {4/x^12};
  \addplot[fmBlue, dashed, thick, domain=0.85:2.6, samples=100] {-4/x^6};
  \addplot[black, very thick, domain=0.92:2.6, samples=150] {4/x^12-4/x^6};
  \draw[dotted] (axis cs:1.1225,0) -- (axis cs:1.1225,-1);
  \draw[dotted] (axis cs:0.8,-1) -- (axis cs:1.1225,-1);
  \node[fmRed, font=\small, anchor=west] at (axis cs:0.95,1.4) {отталкивание $B/r^n$};
  \node[fmBlue, font=\small, anchor=west] at (axis cs:1.5,-0.75) {притяжение $-A/r^m$};
  \draw[dimline] (axis cs:2.2,0) -- node[left, font=\small]{энергия связи} (axis cs:2.2,-1);
  \draw[dashed, fmGray] (axis cs:1.12,-1) -- (axis cs:2.3,-1);
\end{axis}
\end{tikzpicture}
\caption{Энергия взаимодействия двух атомов (потенциал Леннард-Джонса,
$m=6$, $n=12$): сумма притяжения и отталкивания имеет минимум при
равновесном расстоянии $r_0$; глубина минимума~— энергия связи}
\label{fig:LJ}
\end{figure}

\paragraph{Равновесие}
Из условия минимума $\d U/\d r=0$:
\begin{shagi}
\item $\dfrac{\d U}{\d r}=\dfrac{mA}{r^{m+1}}-\dfrac{nB}{r^{n+1}}=0
\;\Rightarrow\; r_0^{\,n-m}=\dfrac{nB}{mA}$, т.\,е. $r_0=\left(\dfrac{nB}{mA}\right)^{1/(n-m)}$.
\item Подставим в $U$: $B/r_0^n=(m/n)\,A/r_0^m$, поэтому
\end{shagi}
\begin{keyeq}\label{eq:U0}
U_0=U(r_0)=-\frac{A}{r_0^m}\left(1-\frac mn\right).
\end{keyeq}
Отталкивание «съедает» долю $m/n$ энергии притяжения: для ионного
кристалла с $m=1$, $n\approx9$~— около 11\,\%.

\begin{smysl}[Что ещё «знает» кривая $U(r)$]
Из одной кривой энергии взаимодействия следуют многие свойства твёрдого
тела:
\begin{itemize}
\item \emph{глубина} минимума~— энергия связи, а с ней теплота
сублимации и (косвенно) температура плавления;
\item \emph{кривизна} в минимуме $\d^2U/\d r^2$~— жёсткость связи, т.\,е.
упругие модули (п.~\ref{sec:elastic}) и частоты колебаний атомов;
\item \emph{асимметрия} ямы (отталкивание круче притяжения)~— тепловое
расширение: при нагреве атомы колеблются с большей амплитудой, и среднее
расстояние растёт (п.~\ref{sec:thermal}).
\end{itemize}
Поэтому температура плавления, модули упругости и коэффициент теплового
расширения так согласованно меняются по периодической системе
(п.~\ref{ssec:atoms}).
\end{smysl}

\paragraph{Энергия кристалла}
Если энергия пары не зависит от других частиц, энергия кристалла из $N$
частиц равна сумме энергий всех пар:
\begin{equation}\label{eq:U-crystal}
U=\frac12\sum_{i=1}^N\sum_{k\ne i}U_{ik},
\end{equation}
где множитель $\frac12$ исключает двойной учёт каждой пары. (Этим же
выражением мы пользовались в п.~2.1.2 при оценке энергии
вакансии.) Для металлов и ковалентных кристаллов приближение парных сил
грубое: энергия связи зависит от окружения атома, и нужны более
сложные модели.
